\chapter{Central limit theorem for martingales}
\label{chap:pre_clt}

We formalize the martingale central limit theorem in Lindeberg form, in dimension
one. This is the analytic core of Lemma~\ref{lem:componentwise}, and through it of
the asymptotic normality Theorem~\ref{thm:normality} and the test calibration
Lemma~\ref{lem:pearson}. A one-dimensional statement suffices: the coordinates of
the estimator error are driven by martingales with vanishing cross-variation
(Lemma~\ref{lem:qv_orthogonal}), so the vector CLT follows coordinatewise via the
Cramér--Wold device or directly from the orthogonality.

\textbf{Formalization note.} Mathlib has the i.i.d.\ one-dimensional CLT
(\texttt{Probability/CentralLimitTheorem.lean}), proved by convergence of
characteristic functions, and the supporting API: \texttt{charFun}
(\texttt{MeasureTheory/Measure/CharacteristicFunction/Basic.lean}), the fact that
\texttt{charFun} determines a measure, and the Gaussian characteristic function
(\texttt{Distributions/Gaussian/CharFun.lean}). What is missing is the martingale
version. The proof reuses the same final step (charFun convergence $\Rightarrow$
convergence in distribution) but replaces independence by a
conditional-characteristic-function argument.

\section{Analytic ingredient}
\label{sec:pre_clt_analytic}

\begin{lemma}[Complex exponential expansion]\label{lem:clt_exp_bound}
For all $x \in \R$,
\[
  \Big| e^{ix} - 1 - ix + \tfrac{x^2}{2} \Big| \le \min\!\Big(x^2,\ \tfrac{|x|^3}{6}\Big).
\]
\end{lemma}

\begin{proof}
Standard: integrate the remainder of the Taylor expansion of $e^{ix}$; the two
bounds come from estimating the remainder integral to first and to second order.
\end{proof}

\begin{lemma}[Conditional characteristic increment]\label{lem:clt_cond_char}
\uses{lem:clt_exp_bound, def:pred_qv}
Let $(M_{n,i})_{i\le k_n}$ be a martingale array with increments $\Delta_{n,i}$.
For fixed $t$,
\[
  \E\big[e^{it\Delta_{n,i}}\mid\filt_{i-1}\big]
  = 1 - \tfrac{t^2}{2}\,\E[\Delta_{n,i}^2\mid\filt_{i-1}] + r_{n,i},
\]
where $\sum_i |r_{n,i}|$ is controlled by the conditional Lindeberg quantity
$\sum_i \E[\Delta_{n,i}^2\indic_{\{|\Delta_{n,i}|>\varepsilon\}}\mid\filt_{i-1}]$
plus $\varepsilon\, t\, \langle M_n\rangle_{k_n}$.
\end{lemma}

\begin{proof}
Take conditional expectation in Lemma~\ref{lem:clt_exp_bound}: the linear term
vanishes since $\E[\Delta_{n,i}\mid\filt_{i-1}]=0$, and the remainder is bounded
by $\E[\min(\Delta_{n,i}^2,|\Delta_{n,i}|^3/6)\mid\filt_{i-1}]$, split at the
threshold $\varepsilon$ into a Lindeberg part and a part $\le \varepsilon t
\E[\Delta_{n,i}^2\mid\filt_{i-1}]$.
\end{proof}

\section{Convergence of characteristic functions}
\label{sec:pre_clt_char}

\begin{lemma}[Telescoping product]\label{lem:clt_product}
\uses{lem:clt_cond_char}
With the notation above, if $\langle M_n\rangle_{k_n} \xrightarrow{\Prob}\sigma^2$
and the conditional Lindeberg condition holds, then
\[
  \E\big[e^{it M_{n,k_n}}\big] \longrightarrow \exp\!\Big(-\tfrac{\sigma^2 t^2}{2}\Big)
  \qquad\text{for every } t\in\R.
\]
\end{lemma}

\begin{proof}
\uses{lem:clt_cond_char, lem:qv_incr}
Write $e^{itM_{n,k_n}} = \prod_i e^{it\Delta_{n,i}}$ and compare with the
predictable product $\prod_i(1 - \tfrac{t^2}{2}\E[\Delta_{n,i}^2\mid\filt_{i-1}])
\approx \exp(-\tfrac{t^2}{2}\langle M_n\rangle_{k_n})$. The telescoping bound uses
$|\prod a_i - \prod b_i| \le \sum_i |a_i - b_i|$ for terms of modulus $\le1$,
with $a_i - b_i = r_{n,i}$ from Lemma~\ref{lem:clt_cond_char}, whose sum vanishes
by Lindeberg. Then $\langle M_n\rangle_{k_n}\xrightarrow{\Prob}\sigma^2$ gives the
limit $\exp(-\sigma^2 t^2/2)$.
\end{proof}

\begin{lemma}[Characteristic functions determine the limit]\label{lem:levy}
The map sending a probability measure to its characteristic function is injective,
and if $\E[e^{itX_n}]\to\exp(-\sigma^2 t^2/2)$ for all $t$, then $X_n$ converges in
distribution to $\mathcal{N}(0,\sigma^2)$.
\end{lemma}

\begin{proof}
Injectivity is \texttt{MeasureTheory.Measure.ext\_of\_charFun} in Mathlib. The
target $\exp(-\sigma^2 t^2/2)$ is the characteristic function of
$\mathcal{N}(0,\sigma^2)$ (\texttt{Distributions/Gaussian/CharFun.lean}). The
implication ``pointwise charFun convergence to a continuous limit charFun
$\Rightarrow$ convergence in distribution'' is Lévy's continuity theorem; Mathlib
has the special case used in its i.i.d.\ CLT
(\texttt{Probability/CentralLimitTheorem.lean}), which is exactly what is needed
here since the limit is a fixed Gaussian.
\end{proof}

\textbf{Formalization note.} A fully general Lévy continuity theorem (arbitrary
tight limits) is not yet in Mathlib, but is not required: we only need
convergence to a fixed, known Gaussian limit, which the CLT file already handles.

\section{The theorem}
\label{sec:pre_clt_thm}

\begin{theorem}[Martingale CLT, Lindeberg form]\label{thm:mart_clt}
\uses{lem:clt_product, lem:levy}
Let $(M_{n,i})_{i \le k_n}$ be a martingale array (with respect to filtrations
$(\filt_{n,i})_i$) with $M_{n,0}=0$, predictable quadratic variation
$\langle M_n\rangle_{k_n}\xrightarrow{\Prob}\sigma^2$, and satisfying the
conditional Lindeberg condition: for every $\varepsilon>0$,
\[
  \sum_{i}\E\big[\Delta_{n,i}^2\indic_{\{|\Delta_{n,i}|>\varepsilon\}}\mid\filt_{n,i-1}\big]
  \xrightarrow{\Prob} 0.
\]
Then $M_{n,k_n}\eqdist\mathcal{N}(0,\sigma^2)$.
\end{theorem}

\begin{proof}
Combine Lemma~\ref{lem:clt_product} (convergence of characteristic functions) with
Lemma~\ref{lem:levy} (the limit is $\mathcal{N}(0,\sigma^2)$).
\end{proof}

\begin{corollary}[Self-normalized componentwise form]\label{cor:mart_clt_componentwise}
\uses{thm:mart_clt, lem:qv_orthogonal}
Let $M_{n,k}=\sum_{i\le n}X_{i,k}(\xi_{i,k}-\theta_k)$ as in
Definition~\ref{def:M}--\ref{def:Q}, with $N_{n,k}\to\infty$ a.s.\ for each $k$.
Then, with $D_n=\mathrm{diag}(\sqrt{N_{n,1}},\dots,\sqrt{N_{n,K}})$,
$D_n(\estParams_n-\Params)\eqdist\mathcal{N}(0,\mathrm{diag}(V_1,\dots,V_K))$.
\end{corollary}

\begin{proof}
\uses{thm:mart_clt, lem:qv_orthogonal}
Apply Theorem~\ref{thm:mart_clt} to the time-changed martingale
$i\mapsto Q_{\tau_k(i),k}$ stopped when $N_{\cdot,k}=i$; its quadratic variation
normalized by $N_{n,k}$ tends to $V_k$, and the conditional Lindeberg condition
holds because the $\xi_{i,k}$ are i.i.d.\ with finite variance (dominated
convergence, cf.\ the proof sketch of Lemma~\ref{lem:componentwise}). Distinct
coordinates have vanishing cross-variation by Lemma~\ref{lem:qv_orthogonal}, so
the joint limit is the stated diagonal Gaussian.
\end{proof}
